\section{Methodology}
\label{sec:method}
\vspace{-1mm}

We present \ours, a framework for EEG emotion recognition that distills relational knowledge from LLM emotion representations into an EEG encoder. Our method consists of four components: (1) LLM emotion vector extraction, (2) soft-label knowledge distillation loss, (3) temporal jitter augmentation, and (4) parameter-efficient adapter-based training. An overview is shown in~\cref{fig:pipeline}.

% === PIPELINE FIGURE (TikZ) ===
\begin{figure*}[!t]
    \centering
    \resizebox{\textwidth}{!}{
    \begin{tikzpicture}[
        node distance=1.2cm and 1.5cm,
        block/.style={rectangle, draw, rounded corners, minimum height=1cm, minimum width=2.2cm, align=center, font=\small},
        frozen/.style={block, fill=blue!8, draw=blue!40},
        trainable/.style={block, fill=green!10, draw=green!50!black},
        loss/.style={block, fill=red!8, draw=red!40, dashed},
        data/.style={block, fill=gray!10, draw=gray!40},
        arrow/.style={-{Stealth[length=2.5mm]}, thick},
        label/.style={font=\scriptsize\itshape, text=gray},
    ]
        % Top branch: EEG pipeline
        \node[data] (eeg) {Raw EEG\\$(B, C, P, 200)$};
        \node[frozen, right=of eeg] (backbone) {CBraMod\\Backbone\\{\scriptsize (frozen, 4.9M)}};
        \node[trainable, right=of backbone] (adapter) {Adapter\\Layers\\{\scriptsize ($r{=}64$, 310K)}};
        \node[trainable, above right=0.5cm and 1.2cm of adapter] (classifier) {Classifier\\{\scriptsize (flatten + MLP)}};
        \node[trainable, below right=0.5cm and 1.2cm of adapter] (proj) {3L-BN\\Projection\\{\scriptsize $(200 \to 512 \to D)$}};
        \node[loss, right=of classifier] (ce) {$\mathcal{L}_{\text{CE}}$};
        \node[loss, right=of proj] (kd) {$\mathcal{L}_{\text{KD}}$};

        % Bottom branch: LLM pipeline
        \node[data, below=2.5cm of eeg] (stories) {Emotion\\Stories};
        \node[frozen, right=of stories] (llm) {LLM\\{\scriptsize (Qwen 0.5B)}\\{\scriptsize frozen}};
        \node[trainable, right=of llm] (vectors) {Emotion\\Vectors\\{\scriptsize $(C, D)$}};
        \node[data, right=of vectors] (sim) {Similarity\\Matrix\\{\scriptsize $\mathbf{S} \in \mathbb{R}^{C \times C}$}};

        % Arrows
        \draw[arrow] (eeg) -- (backbone);
        \draw[arrow] (backbone) -- (adapter);
        \draw[arrow] (adapter) -- (classifier);
        \draw[arrow] (adapter) -- (proj);
        \draw[arrow] (classifier) -- (ce);
        \draw[arrow] (proj) -- (kd);
        \draw[arrow] (stories) -- (llm);
        \draw[arrow] (llm) -- (vectors);
        \draw[arrow] (vectors) -- (sim);
        \draw[arrow] (sim) -- (kd);

        % Labels
        \node[label, above=0.1cm of backbone] {Pretrained on TUEG};
        \node[label, above=0.1cm of adapter] {Task-specific};
        \node[label, below=0.1cm of llm] {67\textsuperscript{th} percentile layer};

        % Final loss
        \node[loss, right=3.5cm of adapter] (total) {$\mathcal{L} = \mathcal{L}_{\text{CE}} + \lambda \mathcal{L}_{\text{KD}}$};
        \draw[arrow, dashed] (ce) -- (total);
        \draw[arrow, dashed] (kd) -- (total);
    \end{tikzpicture}
    }
    \caption{\textbf{Overview of \ours.} The EEG signal is processed by a frozen pretrained backbone (CBraMod or LaBraM) with trainable adapter layers. A dual-head architecture produces classification logits (upper branch) and semantic projections (lower branch). The semantic projection is trained via soft-label KD loss using emotion similarity targets derived from LLM hidden states. The LLM branch (bottom) runs offline to extract emotion vectors.}
    \label{fig:pipeline}
\end{figure*}

\subsection{LLM Emotion Vector Extraction}
\label{sec:llm_vectors}
\vspace{-1mm}

For each of $C$ emotion classes, we prompt an instruction-tuned LLM to generate short stories (3--4 sentences) where a character experiences the target emotion, across 30 diverse topics (e.g., \emph{``a morning commute,''} \emph{``a hospital visit''}). We extract hidden state activations from the 67\textsuperscript{th}-percentile layer (empirically optimal for emotion encoding~\cite{tak2025emotion_interpretability}), average across tokens and stories, and subtract the global mean to obtain \emph{contrast vectors} $\mathbf{v}_c \in \mathbb{R}^D$ for each emotion $c$. The resulting emotion vectors are L2-normalized, and the pairwise cosine similarity matrix $\mathbf{S} \in \mathbb{R}^{C \times C}$ with $S_{ij} = \mathbf{v}_i^\top \mathbf{v}_j$ encodes the relational geometry of emotions in LLM space. Notably, PCA of these vectors recovers the full two-dimensional circumplex~\cite{russell1980circumplex}: PC1 captures valence (positive vs.\ negative) while PC2 captures arousal, placing sadness (low-arousal negative) apart from anger, fear, and disgust (high-arousal negative). The LLM thus internalizes not just which emotions are ``good'' or ``bad,'' but the complete affective geometry established in human psychology.

\noindent\textbf{Universality across LLM families.} We extract vectors from 11 instruction-tuned models in two families: 7 Qwen2.5 sizes (0.5B, 1.5B, 3B, 7B, 14B, 32B, 72B) and 4 Gemma-4 variants (E2B, E4B, 26B-A4B, 31B), as well as several base/multilingual models for control. Pairwise representational similarity (Spearman $\rho$ between RDMs) is high both within Qwen ($\rho = 0.883 \pm 0.067$) and within Gemma ($\rho = 0.904 \pm 0.055$), and---critically---only marginally lower across families ($\rho = 0.818 \pm 0.081$). The strongest cross-family pair (Qwen-0.5B $\leftrightarrow$ Gemma-4 26B-A4B at $\rho = 0.971$) involves models that differ in family, size, and architecture (dense vs.\ MoE). Base models without instruction tuning (Pythia, $r_{\text{PC1-valence}} = 0.30$) and multilingual base models (Bloom, $r = 0.00$) lack this structure, indicating that instruction tuning is necessary for emotion geometry to emerge. For our main paper SOTA we use Qwen2.5-0.5B-Instruct ($D = 896$, 67th-percentile layer), the smallest model with structured emotion vectors. We also validate the marginally-better Qwen-1.5B layer-8 teacher (\cref{sec:llm_scale}) but find that it does not significantly improve downstream EEG accuracy ($+0.4\%$, $p = 0.25$), so we keep the simpler 0.5B configuration.

\noindent\textbf{Token-count requirement.} A practical method-level finding from our cross-dataset analysis is that the KD loss requires a sufficient number of EEG tokens per sample to influence training. With $\leq$60 tokens (e.g., SEED-V's standard 1\,s preprocessing on either backbone), the KD gradient is negligible and the loss has zero effect; we trace this to the small number of patch-level features over which the soft-label gradient can be distributed. Above $\sim$60 tokens, KD becomes effective, with the magnitude of improvement scaling roughly monotonically with token count up to FACED's 321 tokens per sample (\cref{sec:cross_dataset}). The practical implication is that practitioners applying \ours\ to a new EEG dataset should ensure the input granularity yields at least $\sim$60 spatial-temporal tokens per training sample.

\subsection{Soft-Label Knowledge Distillation Loss}
\label{sec:kd_loss}
\vspace{-1mm}

Standard cross-entropy treats all misclassifications equally. Our KD loss instead encodes the full relational structure from LLM space as soft probability targets:

\begin{equation}
    \mathbf{q}_c = \text{softmax}\left(\frac{\mathbf{S}[c, :]}{\tau_t}\right), \quad
    \mathbf{p} = \text{softmax}\left(\frac{\mathbf{z} \cdot \mathbf{P}^\top}{\tau_s}\right)
\end{equation}

where $\mathbf{q}_c$ is the teacher distribution for class $c$ derived from the LLM similarity matrix $\mathbf{S}$ at temperature $\tau_t$, $\mathbf{z}$ is the L2-normalized EEG projection, $\mathbf{P}$ contains the L2-normalized LLM prototypes, and $\tau_s$ is a learnable student temperature. The KD loss is:

\begin{equation}
    \mathcal{L}_{\text{KD}} = \text{KL}(\mathbf{q}_y \| \mathbf{p})
\end{equation}

The total training objective combines standard cross-entropy with the KD term:
\begin{equation}
    \mathcal{L} = \mathcal{L}_{\text{CE}}(\hat{y}, y) + \lambda \, \mathcal{L}_{\text{KD}}(\mathbf{q}_y, \mathbf{p})
\end{equation}

where $\lambda = 0.1$ (see~\cref{sec:kd_ablation} for sensitivity analysis). This formulation differs from standard KD~\cite{hinton2015distilling} in two ways: (1) we use \emph{separate} teacher and student temperatures ($\tau_t = 0.1$, $\tau_s$ learnable), and (2) the teacher signal comes from a cross-modal similarity structure rather than a same-task teacher model.

\noindent\textbf{Motivation from error patterns.} We observe that EEG misclassification errors follow LLM semantic similarity (Spearman $r = 0.38$, $p = 3.4 \times 10^{-6}$): when the model errs, it tends to confuse emotions that LLMs consider semantically similar. This effect is driven entirely by negative emotion confusions ($64.6\%$ of within-category errors, $p = 3.2 \times 10^{-13}$). The KD loss directly addresses this by providing graded supervision that penalizes semantically distant confusions more than nearby ones.

\subsection{Temporal Jitter Augmentation}
\label{sec:augmentation}
\vspace{-1mm}

We apply four EEG-specific augmentations---Gaussian noise, temporal jitter (random shift $\pm 20$ samples), channel dropout, and amplitude scaling---each with probability $p_{\text{aug}} = 0.6$. Through systematic drop-one-out ablation (\cref{sec:aug_ablation}), we find that \textbf{temporal jitter is the single critical augmentation}: removing it drops balanced accuracy by $-3.9\%$ (from $0.609$ to $0.569$), while removing any other augmentation has negligible effect ($\leq 0.4\%$). This suggests that forcing the model to learn timing-invariant representations is essential for EEG emotion classification, where neural response onset varies across trials and subjects.

\subsection{Adapter-Based Frozen Backbone Training}
\label{sec:adapters}
\vspace{-1mm}

We insert lightweight bottleneck adapter modules~\cite{houlsby2019adapters} after each of the 12 transformer layers in the pretrained CBraMod backbone:

\begin{equation}
    \text{Adapter}(\mathbf{x}) = \mathbf{x} + \mathbf{W}_{\text{up}} \cdot \text{GELU}(\mathbf{W}_{\text{down}} \cdot \mathbf{x})
\end{equation}

where $\mathbf{W}_{\text{down}} \in \mathbb{R}^{d \times r}$ and $\mathbf{W}_{\text{up}} \in \mathbb{R}^{r \times d}$ with bottleneck rank $r$, and $\mathbf{W}_{\text{up}}$ is initialized to zero so the adapter starts as identity. The entire backbone (4.9M parameters) is \emph{frozen}; only the adapter parameters (310K for $r = 64$), the classifier head, and the semantic projection head are trained.

This design choice is motivated by our finding that full fine-tuning slightly \emph{degrades} pretrained representations for emotion classification (\cref{sec:adapter_ablation}). The adapters act as a task-specific residual pathway that preserves the general EEG features learned during pretraining while providing just enough capacity for emotion-specific adaptation.
